% Author: Adam | Module: AI Job Interview Simulator | Project: SmartCareer AI
\section{Limites assumées et perspectives}
\label{sec:limites}

Cette section ne cherche pas la formule de circonstance : chaque limite ci-dessous
correspond à un fait observé, mesuré ou écrit dans le code.

\subsection{Limites mesurées ou observées}

\begin{enumerate}[label=\textbf{L\arabic*.},leftmargin=3.2em,itemsep=4pt]
  \item \textbf{Latence incompatible avec une machine sans GPU.} En CPU pur, une
        séance complète (3 questions, 2 évaluations, 1 amélioration) dépasse les
        10~minutes --- 113,6 + 94,3 + 121,5~s par étape sur les mesures du
        tableau~\ref{tab:etapes-cpu}, plus les appels d'évaluation supplémentaires.
        Le cahier des charges visait un poste à 4~Gio de VRAM : le module y tient,
        mais la démonstration n'est confortable que si le modèle est pré-chargé
        et si le GPU n'est pas pris par le Module 1.
  \item \textbf{Pas de streaming, donc pas de progression token par token.}
        \texttt{stream=false} est délibéré (§\ref{sec:modele-raw}) : la sortie est
        un JSON dont aucun préfixe n'est affichable. Le choix a été
        \emph{explicitement reconsidéré} pendant le projet (une voie NDJSON à
        travers route + service + client + Angular a été évaluée puis écartée) ;
        le compromis retenu est une progression par états côté interface, avec
        message temporisé. Conséquence user-visible : entre le clic et la carte,
        l'utilisateur attend 19--120~s sans rien voir bouger.
        \label{sec:no-stream}
  \item \textbf{Libération de VRAM inefficace pendant une génération en vol.}
        Constatée au parcours navigateur (§\ref{sec:perf-vram}) : la requête
        \texttt{keep\_alive: 0} est acceptée mais le modèle reste chargé
        jusqu'à la fin de la génération. Forcer l'annulation couperait une
        réponse en cours de rédaction ; ce n'est pas fait.
  \item \textbf{Contenu \texttt{raw=true} : pas de mémoire entre appels.}
        Questions, évaluation et amélioration sont trois conversations
        indépendantes. Impossible donc de demander « reformule ta question 2 en
        plus court » ou de faire une séance de suivi cohérente ; chaque prompt
        doit ré-expliquer le contexte, ce qui est aussi ce qui les rend longs.
  \item \textbf{RAG aveugle au français hors glossaire.} Mesuré : une requête
        absurde (« qqq www zzz ») marque 0,292 et une requête de plomberie
        0,213 --- au-dessus du seuil de 0,20. Le corpus est anglais, les
        25~termes du glossaire ne couvrent pas le vocabulaire des métiers hors
        ingénierie (§\ref{sec:rag-limites}).
  \item \textbf{Calibration du score non statistique.} Quatre appels réels
        (§\ref{sec:tests-calibration}) sur une question, sans jeu d'or annoté.
        Un biais de longueur est visible dans les données : la meilleure réponse
        est aussi la plus longue, et rien dans le prompt n'interdit de noter la
        longueur. La note affichée doit être lue comme un feedback formatif,
        pas comme une mesure.
  \item \textbf{\texttt{MIN\_QUESTIONS=3} appliqué silencieusement.} Demander 2
        questions en produit 3, côté interface et côté schéma. Cohérent avec le
        barème minimum, mais non expliqué à l'utilisateur.
  \item \textbf{Troncature silencieuse des réponses longues.} 5000~caractères
        envoyés à l'API répondent 200 avec rognage, pas 422
        (§\ref{sec:sec-limites}). Choix assumé, mais un candidat qui écrit un
        roman n'apprend pas que la fin n'a pas été lue.
  \item \textbf{Aucune adaptation à une offre réelle stockée.} Le dépôt du
        Module 1 ne contient aucune table ni colonne d'offre :
        \texttt{get\_job\_offer\_context()} est un point de branchement documenté
        qui renvoie \texttt{None}. L'offre doit être collée par l'utilisateur.
  \item \textbf{Co-tenance GPU non testée en charge.} Le scénario « le Module 1
        génère un CV pendant que le Module 3 évalue » n'a pas été mesuré
        bout en bout avec les deux modèles chargés en même temps ; la
        recommandation \texttt{num\_gpu=0} repose sur l'arithmétique de la
        table~\ref{tab:numgpu}, pas sur une observation directe.
  \item \textbf{Sécurité : dépendance aux mots-clés d'erreur.} La détection de
        contention mémoire cherche \texttt{memory}/\texttt{cuda}/\texttt{oom}/\texttt{vram}
        dans le message d'Ollama (§\ref{sec:sec-erreurs}) ; un changement de
        libellé côté Ollama la rend muette.
  \item \textbf{Persistance limitée à SQLite et au processus unique.} Connexion
        par appel, \texttt{timeout=30} : aucune garantie sous plusieurs workers
        uvicorn, et le module n'a pas été testé ainsi.
\end{enumerate}

\subsection{Perspectives}

Par coût croissant, et en restant dans le projet d'équipe :

\begin{enumerate}[label=\textbf{(\alph*)},leftmargin=*,itemsep=3pt]
  \item \textbf{Signaler plutôt que rogner silencieusement} : une mention
        « réponse tronquée à 4000~caractères » dans la carte d'évaluation, et la
        même pour \texttt{MIN\_QUESTIONS}. Travail d'interface, aucune mesure
        nouvelle nécessaire.
  \item \textbf{Enrichir la base documentaire en français}, ou plonger les
        fragments dans les deux langues : c'est la réponse directe à L5, et le
        corpus étant écrit par le module, l'effort est éditorial.
  \item \textbf{Jeu d'or annoté} : 20 à 30 paires (question, réponse) notées par
        un recruteur, avec corrélation de Spearman entre le score du module et la
        note humaine. Sans cela, la question « est-ce que ce barème note juste ? »
        reste ouverte (L6). Le Module 3 est le seul des deux modules à pouvoir
        produire cet instrument à coût nul pour le modèle.
  \item \textbf{Test de co-tenance réel} (L10) : charger Qwen + LoRA et Mistral
        simultanément, mesurer la VRAM et les échecs, puis fixer le défaut de
        \texttt{num\_gpu} sur la mesure plutôt que sur le calcul.
  \item \textbf{Traçabilité des assertions} : lier chaque \texttt{expected\_point}
        et chaque \texttt{strength} au fragment RAG qui l'a inspiré, pour passer
        d'une puce « Base documentaire » à une citation vérifiable.
  \item \textbf{Fin de séance et séance de suivi} : une mémoire de séance (la
        liste des questions déjà posées, les notes obtenues) est déjà en base ;
        l'exposer dans le prompt permettrait une progression réaliste entre deux
        simulations, ce que L4 interdit aujourd'hui.
  \item \textbf{Adaptation à une offre} (L9) dès que le Module 1 stocke des
        offres : le branchement \texttt{get\_job\_offer\_context} et l'analyse
        d'écart \texttt{skill\_gaps} sont déjà écrits pour ça.
  \item \textbf{Fine-tuning} : le Module 1 a montré la chaîne
        dataset$\to$LoRA$\to$Ollama. Un jeu de paires (réponse faible, feedback,
        réponse STAR) pourrait être généré par le module lui-même et utilisé pour
        affiner le jugé --- mais cela suppose le jeu d'or de (c) pour ne pas
        optimiser un biais.
  \item \textbf{Déploiement} : aujourd'hui le module suppose Ollama et un venv
        local. Une image conteneurisée avec le modèle pré-chargé supprimerait le
        froid de 109~s (L1), au prix d'une distribution de 4~Go.
\end{enumerate}
